
\subsubsection{Gradient Boosting}
El Gradient Boosting es un algoritmo de aprendizaje supervisado perteneciente a la familia de los métodos de ensamble (ensemble methods), específicamente basado en el enfoque de boosting. Su objetivo principal es construir un modelo predictivo fuerte a partir de la combinación de múltiples modelos débiles, generalmente árboles de decisión de baja profundidad.
El objetivo del Gradient Boosting es minimizar una función de pérdida $L(y, F(x))$, donde $y$ representa el valor real y $F(x)$ la predicción del modelo.(\cite{friedman2001})

El modelo se construye de manera aditiva como:

\begin{equation}
	F_m(x) = F_{m-1}(x) + \gamma_m h_m(x)
\end{equation}

donde:
\begin{itemize}
	\item $F_m(x)$ es el modelo en la iteración $m$
	\item $h_m(x)$ es el modelo débil en la iteración $m$
	\item $\gamma_m$ es el coeficiente de ajuste (learning rate)
\end{itemize}

En cada iteración, el modelo se ajusta a los pseudo-residuos, definidos como el gradiente negativo de la función de pérdida:

\begin{equation}
	r_{im} = -\left[ \frac{\partial L(y_i, F(x_i))}{\partial F(x_i)} \right]_{F(x)=F_{m-1}(x)}
\end{equation}

donde:
\begin{itemize}
	\item $r_{im}$ es el pseudo-residuo para la observación $i$ en la iteración $m$
\end{itemize}

El valor óptimo de $\gamma_m$ se obtiene resolviendo:

\begin{equation}
	\gamma_m = \arg\min_{\gamma} \sum_{i=1}^{n} L\left(y_i, F_{m-1}(x_i) + \gamma h_m(x_i)\right)
\end{equation}

Finalmente, el modelo se actualiza iterativamente hasta alcanzar el número de iteraciones deseado o un criterio de convergencia.


\subsubsection{Prophet}

Prophet es un modelo diseñado para realizar pronósticos de series temporales, destacándose por su robustez y flexibilidad para manejar patrones complejos, estacionalidades y datos con valores atípicos. Está pensado para ser intuitivo y fácil de ajustar, permitiendo su uso tanto por expertos como por no especialistas en análisis estadístico. (\cite{hoarau2022time})

Prophet descompone una serie temporal \( y(t) \) en tres componentes principales, sumados a un término de error aleatorio:

\begin{equation}
y(t) = T(t) + S(t) + H(t) + \varepsilon_t
\end{equation}

Donde:
\( T(t) \): Describe la tendencia o evolución a largo plazo de la serie.
\( S(t) \): Captura los patrones estacionales recurrentes.
\( H(t) \): Representa los efectos asociados a eventos específicos o festivos.
\( \varepsilon_t \): Corresponde al ruido aleatorio, que se modela como ruido gaussiano con media cero.

Componente de Tendencia (\( T(t) \))
El modelo permite dos enfoques para capturar tendencias:

1. Crecimiento Lineal: Modela la tendencia como una línea recta con puntos de cambio:
\begin{equation}
T(t) = (k + a(t) \delta)t + (m + a(t)\gamma)
\end{equation}
Donde:
\( k \) es la tasa de crecimiento inicial.
\( m \) es el intercepto inicial.
\( a(t) \) indica si hay puntos de cambio en \( t \).
\( \delta \) y \( \gamma \) ajustan la pendiente e intercepto en dichos puntos.

2. Crecimiento Logístico: Se utiliza cuando el crecimiento está limitado por una capacidad máxima \( C \):
\begin{equation}
T(t) = \frac{C}{1 + \exp(-k(t - m))}
\end{equation}
Aquí:
\( C \) es el valor máximo alcanzable.
\( k \) es la tasa de crecimiento.
\( m \) es el tiempo en el que la serie alcanza la mitad de \( C \).

Componente Estacional (\( S(t) \))
La estacionalidad se representa mediante una combinación de funciones seno y coseno, que permiten capturar patrones periódicos:

\begin{equation}
S(t) = \sum_{n=1}^{N} \left[ a_n \cos\left(\frac{2 \pi n t}{P}\right) + b_n \sin\left(\frac{2 \pi n t}{P}\right) \right]
\end{equation}

Donde:
\( P \) es el periodo del ciclo estacional (por ejemplo, 365 días para patrones anuales).
\( a_n \) y \( b_n \) son coeficientes que determinan la amplitud y la forma del patrón.

Componente de Eventos y Festividades (\( H(t) \)). 
Este componente permite incluir efectos asociados a eventos específicos, modelados como variables indicadoras:

\begin{equation}
H(t) = \sum_{i=1}^{L} \text{indicador}(t \in \text{evento}_i) \cdot \beta_i
\end{equation}

Donde:
\( \text{evento}_i \) representa un evento particular.
\( \beta_i \) es la magnitud del impacto causado por dicho evento.



\subsubsection{XGBoost}
XGBoost (Extreme Gradient Boosting) es un algoritmo de aprendizaje supervisado basado en árboles de decisión y técnicas de boosting. Fue desarrollado como una implementación optimizada del método Gradient Boosting, con mejoras orientadas al rendimiento computacional, regularización del modelo y escalabilidad para grandes volúmenes de datos.
Este algoritmo ha sido ampliamente utilizado en problemas de clasificación, regresión y ranking, debido a su capacidad para obtener alta precisión predictiva y manejar eficientemente datos de gran dimensión. Una de sus principales ventajas consiste en incorporar mecanismos de penalización de complejidad que permiten controlar el sobreajuste y mejorar la capacidad de generalización del modelo. (\cite{chen2016xgboost})

XGBoost construye un modelo de manera secuencial, donde cada nuevo árbol se incorpora con el propósito de corregir los errores generados por los árboles previamente entrenados.

A diferencia del Gradient Boosting tradicional, XGBoost utiliza una expansión de segundo orden de la función de pérdida mediante gradiente y hessiano. Esto permite una optimización más precisa y estable durante el entrenamiento.

El modelo final se expresa como la suma de árboles de decisión:

\begin{equation}
\hat{y}_i = \sum_{k=1}^{K} f_k(x_i), \qquad f_k \in \mathcal{F}
\end{equation}

donde:

\begin{itemize}
\item $\hat{y}_i$ es la predicción para la observación $i$
\item $f_k$ representa el árbol $k$
\item $\mathcal{F}$ es el espacio de árboles de regresión
\item $K$ es el número total de árboles
\end{itemize}


El objetivo de XGBoost es minimizar una función objetivo compuesta por una función de pérdida y un término de regularización:

\begin{equation}
\mathcal{L}(\phi) = \sum_{i=1}^{n} l(y_i,\hat{y}_i) + \sum_{k=1}^{K} \Omega(f_k)
\end{equation}

donde:

\begin{itemize}
\item $l(y_i,\hat{y}_i)$ es la función de pérdida
\item $\Omega(f_k)$ es el término de regularización del árbol
\end{itemize}

La función de regularización está definida como:

\begin{equation}
\Omega(f) = \gamma T + \frac{1}{2}\lambda \sum_{j=1}^{T} w_j^2
\end{equation}

donde:

\begin{itemize}
\item $T$ es el número de hojas del árbol
\item $w_j$ es el peso asociado a la hoja $j$
\item $\gamma$ penaliza el número de hojas
\item $\lambda$ controla la magnitud de los pesos
\end{itemize}

En la iteración $t$, el modelo se actualiza agregando un nuevo árbol:

\begin{equation}
\hat{y}_i^{(t)} = \hat{y}_i^{(t-1)} + f_t(x_i)
\end{equation}

Para facilitar la optimización, XGBoost utiliza una aproximación de segundo orden de la función objetivo:

\begin{equation}
\mathcal{L}^{(t)} \approx \sum_{i=1}^{n}
\left[
l(y_i,\hat{y}_i^{(t-1)})
+ g_i f_t(x_i)
+ \frac{1}{2} h_i f_t^2(x_i)
\right]
+ \Omega(f_t)
\end{equation}

donde:

\begin{equation}
g_i = \frac{\partial l(y_i,\hat{y}_i^{(t-1)})}{\partial \hat{y}_i^{(t-1)}}
\end{equation}

\begin{equation}
h_i = \frac{\partial^2 l(y_i,\hat{y}_i^{(t-1)})}{\partial (\hat{y}_i^{(t-1)})^2}
\end{equation}

Aquí, $g_i$ representa el gradiente de primer orden y $h_i$ el hessiano o derivada de segundo orden.



\subsubsection{Random Forest}

Random Forest es un algoritmo de aprendizaje supervisado basado en técnicas de ensamble, propuesto por Leo Breiman en 2001. Su principio fundamental consiste en construir múltiples árboles de decisión durante el entrenamiento y combinar sus predicciones para obtener un modelo más robusto y con mayor capacidad de generalización.

Este algoritmo es ampliamente utilizado en problemas de clasificación y regresión debido a su capacidad para reducir el sobreajuste que suele presentarse en árboles de decisión individuales. Además, Random Forest presenta buen desempeño frente a datos con alta dimensionalidad, ruido y relaciones no lineales entre variables.(\cite{breiman2001randomforest})


Random Forest pertenece a la familia de métodos \textit{bagging} (\textit{bootstrap aggregating}). El procedimiento consiste en generar múltiples subconjuntos de entrenamiento mediante muestreo aleatorio con reemplazo a partir del conjunto de datos original.

Para cada subconjunto se entrena un árbol de decisión independiente. Durante la construcción de cada nodo, el algoritmo selecciona aleatoriamente un subconjunto de variables predictoras, lo cual introduce diversidad entre los árboles.

En problemas de clasificación, la predicción final se obtiene mediante votación mayoritaria entre los árboles. En problemas de regresión, la salida final corresponde al promedio de las predicciones individuales.


Sea un conjunto de entrenamiento:

\begin{equation}
	D = \{(x_i, y_i)\}_{i=1}^{n}
\end{equation}

donde $x_i$ representa el vector de variables de entrada y $y_i$ la variable objetivo.

Se generan $B$ muestras bootstrap del conjunto de entrenamiento, y para cada muestra se entrena un árbol de decisión $h_b(x)$.

En problemas de clasificación, la predicción final se define como:

\begin{equation}
	\hat{y} = \operatorname{mode}\{h_1(x), h_2(x), \dots, h_B(x)\}
\end{equation}

donde $\operatorname{mode}$ representa la clase más frecuentemente predicha por los árboles.

En problemas de regresión, la predicción final se obtiene mediante el promedio:

\begin{equation}
	\hat{y} = \frac{1}{B}\sum_{b=1}^{B} h_b(x)
\end{equation}

Durante la construcción de cada árbol, en cada nodo se selecciona aleatoriamente un subconjunto de $m$ variables entre las $p$ variables disponibles, con la restricción:

\begin{equation}
	m < p
\end{equation}

Este proceso reduce la correlación entre los árboles y mejora la capacidad de generalización del modelo.

\subsubsection{K-Nearest Neighbors}

K-Nearest Neighbors (KNN) es un algoritmo de aprendizaje supervisado utilizado en problemas de clasificación y regresión. Fue propuesto como un método no paramétrico basado en la similitud entre observaciones, donde la predicción de una nueva instancia se realiza a partir de los ejemplos más cercanos presentes en el conjunto de entrenamiento.Una de las principales características de KNN es que no construye explícitamente un modelo durante la fase de entrenamiento. En lugar de ello, almacena los datos observados y realiza el proceso de inferencia cuando se presenta una nueva observación. Debido a su simplicidad conceptual y facilidad de implementación, es uno de los algoritmos fundamentales dentro del aprendizaje automático.

El principio de funcionamiento de KNN consiste en identificar las $k$ observaciones del conjunto de entrenamiento más cercanas a una nueva observación, de acuerdo con una medida de distancia definida previamente.
En problemas de clasificación, la clase asignada corresponde a la categoría mayoritaria entre los vecinos más próximos. En problemas de regresión, la predicción se obtiene como el promedio de los valores de dichos vecinos.
El desempeño del algoritmo depende principalmente de tres factores: el valor de $k$, la métrica de distancia utilizada y la escala de las variables predictoras. (\cite{hastie2009elementsKNN})


Sea un conjunto de entrenamiento definido como:

\begin{equation}
	D = \{(x_i, y_i)\}_{i=1}^{n}
\end{equation}

donde $x_i \in \mathbb{R}^{p}$ representa el vector de características y $y_i$ la variable objetivo asociada.

Dada una nueva observación $x$, se calcula la distancia entre $x$ y cada observación del conjunto de entrenamiento. Una de las métricas más utilizadas es la distancia euclidiana:

\begin{equation}
	d(x, x_i) = \sqrt{\sum_{j=1}^{p}(x_j - x_{ij})^2}
\end{equation}

Posteriormente, se seleccionan los $k$ vecinos más cercanos:

\begin{equation}
	N_k(x) = \{x_{(1)}, x_{(2)}, \dots, x_{(k)}\}
\end{equation}

donde $x_{(j)}$ representa la observación correspondiente al $j$-ésimo vecino más próximo.

En problemas de clasificación, la predicción se obtiene mediante votación mayoritaria:

\begin{equation}
	\hat{y} = \operatorname{mode}\{y_i : x_i \in N_k(x)\}
\end{equation}

En problemas de regresión, la predicción se define como el promedio de las respuestas de los vecinos seleccionados:

\begin{equation}
	\hat{y} = \frac{1}{k} \sum_{x_i \in N_k(x)} y_i
\end{equation}

El valor de $k$ controla el grado de suavizado del modelo. Valores pequeños pueden generar alta sensibilidad al ruido, mientras que valores grandes tienden a producir fronteras de decisión más estables.



\subsubsection{Holt-Winters}

El método de Holt-Winters es una técnica de suavizamiento exponencial utilizada para el análisis y pronóstico de series temporales que presentan tendencia y estacionalidad. Fue desarrollado como una extensión del método de suavizamiento exponencial de Holt, incorporando un componente adicional que permite modelar patrones estacionales recurrentes.

Este método es ampliamente utilizado en problemas de pronóstico de demanda, ventas, consumo energético y otras aplicaciones donde los datos presentan comportamiento temporal estructurado. Su principal ventaja radica en su capacidad para capturar simultáneamente nivel, tendencia y estacionalidad dentro de una misma formulación.(\cite{hyndman2018forecastingHW})


El método de Holt-Winters descompone una serie temporal en tres componentes fundamentales:

\begin{itemize}
	\item nivel (\(l_t\))
	\item tendencia (\(b_t\))
	\item estacionalidad (\(s_t\))
\end{itemize}

En cada instante de tiempo, estos componentes se actualizan utilizando ponderaciones exponenciales, otorgando mayor importancia a las observaciones recientes.

Existen dos variantes principales del método: aditiva y multiplicativa. La versión aditiva se utiliza cuando la magnitud de la estacionalidad permanece aproximadamente constante en el tiempo, mientras que la versión multiplicativa es apropiada cuando la amplitud estacional varía proporcionalmente con el nivel de la serie.


Sea \(y_t\) una serie temporal observada en el instante \(t\), y sea \(m\) la longitud del período estacional.

En la versión aditiva, las ecuaciones de actualización son las siguientes.

Actualización del nivel:

\begin{equation}
	l_t = \alpha (y_t - s_{t-m}) + (1-\alpha)(l_{t-1} + b_{t-1})
\end{equation}

Actualización de la tendencia:

\begin{equation}
	b_t = \beta (l_t - l_{t-1}) + (1-\beta)b_{t-1}
\end{equation}

Actualización de la estacionalidad:

\begin{equation}
	s_t = \gamma (y_t - l_t) + (1-\gamma)s_{t-m}
\end{equation}

donde:

\begin{itemize}
	\item \(l_t\) representa el nivel de la serie en el instante \(t\)
	\item \(b_t\) representa la tendencia
	\item \(s_t\) representa el componente estacional
	\item \(\alpha\), \(\beta\) y \(\gamma\) son parámetros de suavizamiento tales que \(0 < \alpha, \beta, \gamma < 1\)
\end{itemize}

El pronóstico para \(h\) períodos hacia adelante se calcula mediante:

\begin{equation}
	\hat{y}_{t+h} = l_t + hb_t + s_{t-m+h_m}
\end{equation}

donde:

\begin{equation}
	h_m = ((h-1) \bmod m) + 1
\end{equation}

En la variante multiplicativa, el componente estacional interviene de forma proporcional.

Actualización del nivel:

\begin{equation}
	l_t = \alpha \left(\frac{y_t}{s_{t-m}}\right) + (1-\alpha)(l_{t-1}+b_{t-1})
\end{equation}

Actualización de la estacionalidad:

\begin{equation}
	s_t = \gamma \left(\frac{y_t}{l_t}\right) + (1-\gamma)s_{t-m}
\end{equation}

Pronóstico:

\begin{equation}
	\hat{y}_{t+h} = (l_t + hb_t)\, s_{t-m+h_m}
\end{equation}



\subsubsection{Naive}

El método Naive es una técnica básica de pronóstico de series temporales utilizada como referencia en modelos de predicción. Su principio consiste en asumir que el valor futuro de una variable será igual al valor observado en el período inmediatamente anterior.

A pesar de su simplicidad, el método Naive posee gran importancia práctica, ya que suele emplearse como modelo base de comparación para evaluar el desempeño de métodos de pronóstico más complejos. En numerosas aplicaciones, especialmente cuando la serie presenta alta persistencia temporal, puede producir resultados competitivos.(\cite{hyndman2018forecastingNV})


El enfoque del método Naive parte del supuesto de que la mejor estimación del siguiente valor observado es el último valor disponible de la serie.

Este procedimiento no requiere estimación de parámetros ni entrenamiento de modelos. La información histórica se resume únicamente en la observación más reciente.

En series temporales con estacionalidad también existe una variante denominada \textit{Seasonal Naive}, en la cual el valor futuro se estima utilizando el valor correspondiente al mismo período estacional anterior.


Sea \(y_t\) una serie temporal observada en el instante \(t\).

El pronóstico de un período hacia adelante se define como:

\begin{equation}
	\hat{y}_{t+1} = y_t
\end{equation}

De forma general, el pronóstico para \(h\) períodos hacia adelante está dado por:

\begin{equation}
	\hat{y}_{t+h} = y_t
\end{equation}

Esto implica que todas las predicciones futuras corresponden al último valor observado de la serie.

En el caso de series temporales estacionales, sea \(m\) la longitud del período estacional. El método Seasonal Naive se define como:

\begin{equation}
	\hat{y}_{t+h} = y_{t+h-m}
\end{equation}

donde el pronóstico se obtiene utilizando el valor observado en el mismo período de la estación anterior.

Este método es especialmente útil como línea base de comparación, ya que permite determinar si modelos más sofisticados aportan mejoras reales en capacidad predictiva.

\subsection{Computación en la nube}


La computación en la nube (\textit{cloud computing}) es un modelo de provisión de recursos informáticos bajo demanda a través de internet. Permite acceder a capacidad de procesamiento, almacenamiento, redes y servicios de software sin necesidad de que la organización mantenga infraestructura física propia.

De acuerdo con el National Institute of Standards and Technology (NIST), este paradigma se caracteriza por acceso ubicuo a recursos compartidos, aprovisionamiento rápido, elasticidad y escalabilidad. Estas propiedades permiten que los sistemas se adapten dinámicamente a variaciones en carga de trabajo, volumen de datos y número de usuarios.

En proyectos de analítica de datos y pronóstico, la computación en la nube facilita el despliegue de aplicaciones que integran procesamiento, visualización y acceso remoto a resultados. Además, favorece la reproducibilidad, la disponibilidad y la colaboración entre usuarios, lo cual resulta especialmente relevante en entornos de investigación aplicada.

\subsection{Streamlit}

Streamlit es un framework de código abierto orientado al desarrollo de aplicaciones web interactivas para ciencia de datos y aprendizaje automático utilizando el lenguaje Python.

Su principal característica consiste en permitir la construcción de interfaces de usuario de manera declarativa y con bajo nivel de complejidad técnica. A partir de scripts en Python, es posible integrar visualizaciones, tablas, métricas, formularios y componentes interactivos que facilitan la exploración de datos y la presentación de resultados analíticos.

En el contexto de una investigación orientada a pronóstico, Streamlit permite materializar los modelos desarrollados en una aplicación accesible mediante navegador web, facilitando la interacción con usuarios finales y la consulta de resultados en tiempo real.

\subsection{Gestión de inventarios}

La gestión de inventarios comprende el conjunto de actividades orientadas a planificar, controlar y administrar los niveles de existencias dentro de una organización. Su finalidad consiste en asegurar la disponibilidad de productos para atender la demanda, minimizando simultáneamente costos asociados al almacenamiento, reposición, faltantes y pérdidas.(\cite{chopra2019supply})

En empresas comerciales, una adecuada gestión de inventarios permite mantener equilibrio entre el nivel de servicio al cliente y la eficiencia operativa. Cuando los niveles de inventario son insuficientes pueden producirse quiebres de stock y pérdida de ventas; por el contrario, cuando los niveles son excesivos se incrementan los costos de almacenamiento, inmovilización de capital y riesgo de obsolescencia o deterioro.

\subsubsection{Abastecimiento}

El abastecimiento es el proceso mediante el cual una organización planifica, adquiere y repone los bienes necesarios para garantizar la continuidad de sus operaciones. Este proceso comprende actividades como la planificación de requerimientos, selección de proveedores, programación de pedidos y seguimiento de entregas.

En el contexto comercial, las decisiones de abastecimiento dependen en gran medida de la estimación de la demanda futura. Un proceso de abastecimiento basado en información predictiva permite reducir incertidumbre, mejorar la disponibilidad de productos y optimizar la asignación de recursos destinados a compras.

\subsubsection{Productos perecibles}

Los productos perecibles son bienes que presentan una vida útil limitada y cuya calidad se deteriora progresivamente con el transcurso del tiempo. Debido a esta característica, su almacenamiento prolongado incrementa la probabilidad de pérdidas económicas asociadas a vencimiento, deterioro físico o reducción de valor comercial.

La gestión de productos perecibles requiere especial atención en la planificación de inventarios, ya que las decisiones de reposición deben considerar simultáneamente la demanda esperada, la frecuencia de abastecimiento y el horizonte de conservación del producto.

En empresas dedicadas a la comercialización de productos perecibles, un pronóstico adecuado de la demanda permite reducir desperdicios, mejorar la disponibilidad y contribuir a una operación más eficiente.


\subsubsection{Rotación de inventarios}

La rotación de inventarios es un indicador que mide la frecuencia con que los productos son vendidos y repuestos durante un período determinado. Este indicador permite evaluar la velocidad de movimiento de los inventarios y el grado de eficiencia en la utilización de los recursos almacenados.\cite{heizer2017operations}

Matemáticamente, la rotación de inventarios puede expresarse como:

\begin{equation}
	\text{Rotación de inventarios} =
	\frac{\text{Costo de ventas}}{\text{Inventario promedio}}
\end{equation}

Un valor alto de rotación generalmente indica que los productos tienen mayor dinamismo comercial y menor permanencia en almacén. En productos perecibles, una adecuada rotación contribuye a disminuir riesgos de vencimiento y deterioro.
